\documentclass[11pt]{article}
\usepackage[T1]{fontenc}
\usepackage[utf8]{inputenc}
\usepackage[english]{babel}
\usepackage{newpxtext,newpxmath}
\usepackage[a4paper,top=1.6cm,bottom=1.55cm,left=2.1cm,right=2.1cm]{geometry}
\usepackage{microtype}
\usepackage{xcolor}
\usepackage{titlesec}
\usepackage[most]{tcolorbox}
\usepackage{fancyhdr}
\usepackage{enumitem}
\usepackage{booktabs}
\usepackage{listings}
\lstdefinestyle{jsonic}{basicstyle=\ttfamily\footnotesize,backgroundcolor=\color{boxgray},%
  breaklines=true,columns=fullflexible,keepspaces=true,xleftmargin=6pt,%
  framexleftmargin=6pt,aboveskip=4pt,belowskip=4pt}
\lstset{style=jsonic}

\definecolor{efenavy}{HTML}{1F3A5F}
\definecolor{efeteal}{HTML}{0E7A6E}
\definecolor{efeamber}{HTML}{B25A1E}
\definecolor{efeplum}{HTML}{7A4EA8}
\definecolor{efegray}{HTML}{54606E}
\definecolor{efesolve}{HTML}{2C6E49}
\definecolor{boxgray}{HTML}{F3F5F8}
\definecolor{boxamber}{HTML}{FBF3EA}
\definecolor{boxedge}{HTML}{D9E0EA}

\titleformat{\section}{\color{efenavy}\large\bfseries}{}{0pt}{}
\titlespacing{\section}{0pt}{7pt}{3pt}
\titleformat{\subsection}[runin]{\color{efenavy}\bfseries}{}{0pt}{}[:\;]
\titlespacing{\subsection}{0pt}{5pt}{0.5em}

\setlength{\parskip}{0.6pt}
\setlength{\parindent}{0pt}
\setlength{\emergencystretch}{3em}

\pagestyle{fancy}\fancyhf{}
\renewcommand{\headrulewidth}{0.3pt}
\lhead{\small\color{efegray}Scella --- the Threshold \& the Register}
\rhead{}
\cfoot{\small\color{efegray}\thepage}

\newcommand{\mesure}[1]{{\color{efeteal}\textbf{Measured.} }#1\par}
\newcommand{\concept}[1]{{\color{efenavy}\textbf{Concept:} }#1\par}
\newcommand{\exemple}[1]{{\color{efeamber}\textbf{Example:} }#1\par}
\newcommand{\utilite}[1]{{\color{efeplum}\textbf{Real-world use:} }#1\par}
\newcommand{\solves}[2]{{\color{efesolve}\textbf{Solves #1:} }#2\par}

\newtcolorbox{these}{colback=boxgray,colframe=boxedge,boxrule=0.5pt,arc=2.5pt,
  left=8pt,right=8pt,top=6pt,bottom=6pt}
\newtcolorbox{cadre}{colback=boxamber,colframe=efeamber!55,boxrule=0.6pt,arc=2.5pt,
  left=8pt,right=8pt,top=6pt,bottom=6pt}

\begin{document}

\begin{center}
{\Huge\bfseries\color{efenavy} Scella}\\[3pt]
{\large\color{efegray}\itshape the threshold that seals, the register that answers}
\end{center}
\vspace{6pt}

\begin{these}
\textbf{The problem.} A regulated organisation wants to use AI without letting its data leave, and
must later answer a regulator, an auditor or a court \emph{in the language of the text that binds
it} --- without handing over its systems, and without anyone having to take the operator's word for
what happened.

\textbf{What Scella concretely solves:}
\begin{itemize}[leftmargin=1.3em,itemsep=1.5pt,topsep=2pt,parsep=0pt,label=\textendash]
\item keep the \textbf{data and the AI inside your perimeter}, and seal every call that crosses it;
\item \textbf{answer a regulator} in the vocabulary of the text --- AI Act, GDPR, CSDDD, LkSG ---
rather than in the vocabulary of the software;
\item \textbf{survive an incident}: a rollout, a restore, a key rotation, without losing the proof;
\item \textbf{never degrade silently}: what is refused is refused \emph{by name}, with its cause;
\item let a \textbf{third party re-derive the proof offline}, on his own machine, with no access to
your systems and no trust in you.
\end{itemize}
\vspace{2pt}
\textbf{Two products, one receipt.} \textbf{The Threshold} (C3, 20 crates) is the sovereign gateway:
it runs where the data is and seals each call. \textbf{The Register} (C1, 8 crates) turns those seals
into the evidence a \emph{named} regulation asks for. They share the same sealed node, so the
compliance document and the proof cannot drift apart --- and they ship together or apart.
\end{these}

\section{I \;\textbar\; The Threshold (C3): the sealed perimeter}

\subsection{What it is}
A gateway that stands at the edge of your systems. Calls to models pass through it; each one becomes
a \textbf{sealed node} in an append-only chain, and the chain re-proves offline. It starts with one
command and it does not need the internet to prove anything.
\concept{Not a proxy that logs. A threshold that \emph{seals}: what crosses it leaves a mark that
cannot be edited afterwards without breaking the next mark.}
\utilite{Bank, hospital, public administration: the model can be remote, the \emph{evidence} stays
in the building, and the auditor is handed a file rather than an access.}
\solves{(1)}{the AI comes to the data, not the reverse.}

\subsection{Three states, and never a silent success}
Every capability the gateway offers is in exactly one of three states: \emph{served}, \emph{refused
with its cause named}, or \emph{degraded and said so}. The third is the one that costs discipline and
buys trust: a provenance that cannot persist reports \texttt{degraded} \emph{with a reason} in the
response, it does not quietly return 200 and drop the record.
\exemple{Ask for a capability that is switched off and you get a refusal that \emph{names the
switch}. Ask for one whose store is unreachable and you get your answer \emph{plus} the sentence
saying the receipt was not persisted, and why.}
\mesure{The boot guard refuses to start the paying frontier without its preconditions and exits
\textbf{78} \emph{naming the variable} it wanted --- \texttt{EFE\_NANO\_AUDIT\_HMAC\_MASTER},
\texttt{EFE\_OIDC\_PUBLIC\_KEY\_PEM}, \texttt{DATABASE\_URL}. Verified in one command:
\texttt{bash scripts/p1-check.sh}.}
\solves{(4)}{you always know what the system did \emph{not} do, and why.}

\subsection{One authority analyser, zero exemptions}
A single component decides which host the HTTP client will \emph{really} contact --- and it is the
only one allowed to decide. This is not architectural taste: before it existed,
\texttt{127.0.0.1:8080@evil.com} passed as loopback and \textbf{cut TLS towards
\texttt{evil.com}}, on \textbf{875} spelling variants of the same trick.
\concept{A perimeter with two doors has no perimeter. The exemption list is empty, and a test in the
repository fails if anyone adds a second analyser.}
\solves{(1)}{one door, one guard, and a test that notices when a second door appears.}

\subsection{The anchor, and the one failure re-proof is blind to}
Re-proof catches editing. It cannot catch \emph{truncation}: a valid chain's prefix is a valid
chain, so a receipt with its last three steps removed re-proves \textbf{perfectly}. The root and the
node count are therefore published to an append-only log \emph{before} the receipt is served.
\mesure{The truncation is then exposed \textbf{by counting} --- the anchor says three nodes, the
bundle contains two --- \textbf{with no key and no cryptography at all}. Anyone can redo it with
\texttt{grep -c} and \texttt{awk}. Verified end-to-end by \texttt{docker compose up}, which brings
the gateway up, serves a run, exports the bundle, truncates a copy and shows the contradiction.}
\solves{(0)}{the failure mode that re-proof alone will never see.}

\subsection{Surviving a rollout}
A receipt is worth what it is worth \emph{after} the incident. The gateway is started against a real
PostgreSQL, serves a run, is killed; \emph{another process} takes the same port; the receipt is
re-proved from the new process.
\mesure{\texttt{verified $\cdot$ root\_matches $\cdot$ chain\_intact $\cdot$ content\_intact}.
This is the acceptance criterion of the production profile, and it is a script:
\texttt{bash scripts/p1-check.sh} (it starts its own PostgreSQL 16 in a container).}
\solves{(3)}{a restart, a deploy or a crash does not cost you the proof.}

\section{II \;\textbar\; The Register (C1): the answer a regulation asks for}

\subsection{Four texts, by name}
The register produces evidence for \textbf{EU AI Act} (Regulation (EU) 2024/1689, art.~12
record-keeping), \textbf{GDPR} (Regulation (EU) 2016/679), \textbf{CSDDD} (Directive (EU) 2024/1760)
and \textbf{LkSG} (Germany). Each artefact is written in the vocabulary of its text.
\concept{A compliance officer should not have to translate ``sealed DAG node'' into ``record of the
functioning of the system''. The register does the translation, and the mapping is versioned so that
re-stamping an unchanged reading is impossible.}
\utilite{Supply-chain due diligence, high-risk AI record-keeping, processing register: the same
sealed material, three different readings.}
\solves{(2)}{an answer that can be handed over as it is.}

\subsection{Derived, never re-declared}
The evidence is \textbf{derived from the sealed chain} --- not written alongside it. There is no
second source of truth to keep in sync, and no moment at which a human retypes what the system did.
\concept{If the register were declared separately, it could drift from the receipt, and the drift
would be invisible. Deriving makes drift \emph{impossible}: rewriting one breaks the re-proof of the
other.}
\solves{(0) \& (2)}{the compliance document \emph{is} the proof, not a story about it.}

\subsection{Erasure without breaking the chain}
GDPR erasure and record-keeping duties pull in opposite directions. The register resolves it by
\textbf{crypto-shredding}: erasable fields are sealed \emph{encrypted} under a per-case key, and
erasing destroys that key. Not one byte of the sealed graph moves; re-proof still passes.
\concept{An erased field reads \textbf{``destroyed on \emph{date}''} --- never ``unknown''. The two
must never be confused: the first is a dated fact, the second an absence. Telling a regulator that a
case he saw pass never existed is the failure this distinction exists to prevent.}
\mesure{The keyring is durable across restarts, and a destruction that cannot be made durable is
\emph{refused} rather than performed in memory --- a shredding that the next restart would undo is a
lie about a destruction.}
\solves{(2) \& (3)}{the right to erasure and the duty to keep records, in the same artefact.}

\subsection{Joined, not welded}
The Threshold reaches the Register through a single feature, \textbf{off by default} and guarded by
a manifest check. Enabling it is a deployment decision, not a rebuild.
\utilite{Buy the perimeter first because an incident is expensive; add the register when the first
audit is announced. The second purchase does not disturb the first.}
\solves{(2)}{two products that ship together or apart, without a fork.}

\section{III \;\textbar\; What binds them: one receipt}

\subsection{The sealed node}
Both products write the same thing: an \textbf{HMAC chain} of nodes, keyed per tenant, where each
seal covers the canonical encoding of its node \emph{and} the previous chain value. Tampering with
one link makes every following link fail.
\concept{No new cryptography, no home-made construction: SHA-256 (FIPS 180-4) and HMAC (RFC 2104),
and a key derivation whose context string is versioned so that changing it changes every key
deliberately rather than by accident.}

\subsection{A verifier with an empty dependency tree}
The receipt is re-proved by a program whose \textbf{declared dependency tree contains zero crates}.
A third party recompiles it from source in seconds, reads it in an afternoon, and checks your receipt
on his own machine.
\concept{``Trust our verifier'' is not an argument. ``Recompile ours, or write your own from a
20-line specification'' is one --- and the empty dependency tree is what makes the second sentence
true rather than rhetorical.}
\solves{(0)}{offline re-derivation, without your systems and without trusting you.}

\subsection{Conservation: backup, restore, key rotation}
A register meant to last ten years must survive a disk loss and an administrator's departure. Backup
writes an archive that carries what a restore needs to \emph{affirm} something --- and nothing more:
never the key, and only the \emph{fingerprint} of the case keyring, not its contents.
\concept{Key rotation is \textbf{additive}, never substitutive. A \emph{retired} generation stops
sealing new records but still re-proves everything it sealed; only a \emph{destroyed} one stops, and
its receipts then read ``destroyed on \emph{date}'' --- never ``invalid'', which would accuse someone
of tampering.}
\mesure{A restore distinguishes \textbf{fourteen} distinct outcomes --- truncated, extended,
re-sealed, wrong key, destroyed generation, keyring mismatch --- rather than one ``verification
failed''. Among them, one is \emph{proved} rather than guessed: if another known generation re-proves
the archive, the manifest is wrong, and the tool says which generation is right.}
\solves{(3)}{the incident, the ten-year horizon, and the administrator who leaves.}

\begin{cadre}
\textbf{The exact guarantee --- tamper-evident to whoever holds the key, still \emph{not}
trustless.} The anchor log holder
is the key holder, so publishing the root removes \emph{retroactive} rewriting --- it does not remove
trust in the operator. Third-party time-stamping (RFC 3161) is wired: the token is obtained and
deposited verbatim, and \textbf{its signature is verified by the third party} with
\texttt{openssl ts -verify}. Re-implementing that verification inside the repository would require a
full ASN.1/X.509 stack; \emph{no home-made crypto, no third stack} is a standing rule here, and it
protects more than it costs. An asset of \textbf{verifiable discipline}, not a distributed vault.
\end{cadre}

\section{Results --- what it yields, by part}
What the measurements establish, positively, product by product (measured 2026-08-27):
\begin{itemize}[leftmargin=1.3em,itemsep=2.5pt,topsep=3pt,label=\textendash]
\item \textbf{Engineering} --- \textbf{22 crates} across both products (20 for the Threshold, 8 for
the Register, overlapping on the receipt), \textbf{2{,}539 green tests, 0 red}.
\item \textbf{Executable claims} --- \textbf{15 checks} under \texttt{scripts/}. Every statement in
this document is re-runnable in one command; the two that matter most are the rollout
(\texttt{p1-check.sh}, with a real PostgreSQL) and the single command
(\texttt{docker compose up}), both verified end-to-end.
\item \textbf{The Threshold} --- serves without the internet, refuses by name, degrades out loud;
one authority analyser with an \emph{empty} exemption list, proven necessary against 875 spelling
variants of a loopback bypass.
\item \textbf{The Register} --- four named texts; evidence derived from the chain rather than
re-declared; erasure that destroys a key without moving a byte of the graph, and that reports a
\emph{date} rather than an absence.
\item \textbf{The receipt} --- re-proved by a verifier with \textbf{zero declared dependencies};
truncation caught by counting, \textbf{with no key at all}; backup, restore and rotation with a
fourteen-way diagnosis and a retired generation that keeps re-proving.
\item \textbf{The method} --- a re-provable instrument: frozen constants, deterministic results, and
claims that are refuted in writing when measurement contradicts them.
\end{itemize}


% La note de (0)-(4) glissait seule sur une 4e page. `\enlargethispage` est l'outil prevu :
% on etire CETTE page plutot que de retoucher la mise en page de tout le document.
\enlargethispage{4\baselineskip}
\vspace{4pt}
{\color{efegray}\hrule height 0.4pt}
\vspace{2.5pt}
{\footnotesize\color{efegray}(0)--(4) refer to the problems above: (0)~let a third party re-derive the proof offline, without your systems and without trusting you; (1)~keep the data and the AI inside your perimeter; (2)~answer a regulator in the language of the text that binds you; (3)~survive an incident --- rollout, restore, key rotation; (4)~never a silent degradation: what is refused is named.}

\end{document}
